%%=============================================================================
%% Voorwoord
%%=============================================================================

\chapter*{\IfLanguageName{dutch}{Woord vooraf}{Preface}}%
\label{ch:voorwoord}

%% TODO:
%% Het voorwoord is het enige deel van de bachelorproef waar je vanuit je
%% eigen standpunt (``ik-vorm'') mag schrijven. Je kan hier bv. motiveren
%% waarom jij het onderwerp wil bespreken.
%% Vergeet ook niet te bedanken wie je geholpen/gesteund/... heeft

De keuze van dit onderwerp voor mijn bachelorproef komt voort uit met mijn interesse in cybersecurity en omdat ik graag mensen informeer en iets bijleer. In het nieuws komen steeds meer gevallen van phishing en hacking voor en van daaruit groeit het belang van het beschermen van digitale gegevens. Dit maakt het onderwerp over het gebruik van een wachtwoordmanager extra interessant omdat het blijkt dat, ondanks de beschikbaarheid, weinig mensen gekent zijn met het gebruik van deze tool. Uit de literatuurstudie heb ik kunnen achterhalen wat de drempels waren voor gebruikers in het niet gebruiken van een wachtwoordmanager en hoe ik deze mensen beter kan ondersteunen.
Ik heb tijdens deze bachelorproef niet alleen het menselijke aspect bestudeerd maar ook het technische, wat een heel leerrijk proces was. 

Graag wil ik mijn promotor, meneer Bosteels en mijn co-promotor, meneer Clauwaert, bedanken voor de frequente begeleiding, de constructieve feedback en ondersteuning. Verder wil ook mijn vriendin Laura, mijn vrienden en familie bedanken voor de mentale steun en motivatie tijdens dit proces.
